% !TeX program = pdflatex
\documentclass[11pt,a4paper]{amsart}

\newcommand{\notelanguage}{finnish}
% Shared preamble for course notes and exercise sets.

\usepackage[T1]{fontenc}
\usepackage{lmodern}
\providecommand{\notelanguage}{english}
\usepackage[finnish,english]{babel}
\AtBeginDocument{\expandafter\selectlanguage\expandafter{\notelanguage}}

\newcommand{\notesfinnishlanguage}{finnish}
\ifx\notelanguage\notesfinnishlanguage
\newcommand{\notestheoremname}{Lause}
\newcommand{\noteslemmaname}{Lemma}
\newcommand{\notespropositionname}{Propositio}
\newcommand{\notescorollaryname}{Seurauslause}
\newcommand{\notesdefinitionname}{Määritelmä}
\newcommand{\notesexamplename}{Esimerkki}
\newcommand{\notesproblemname}{Tehtävä}
\newcommand{\notesexercisename}{Tehtävä}
\newcommand{\notesremarkname}{Huomautus}
\newcommand{\notessolutionname}{Ratkaisu}
\else
\newcommand{\notestheoremname}{Theorem}
\newcommand{\noteslemmaname}{Lemma}
\newcommand{\notespropositionname}{Proposition}
\newcommand{\notescorollaryname}{Corollary}
\newcommand{\notesdefinitionname}{Definition}
\newcommand{\notesexamplename}{Example}
\newcommand{\notesproblemname}{Problem}
\newcommand{\notesexercisename}{Exercise}
\newcommand{\notesremarkname}{Remark}
\newcommand{\notessolutionname}{Solution}
\fi

\usepackage[a4paper,margin=30mm]{geometry}
\usepackage{microtype}
\usepackage{graphicx}
\usepackage{enumitem}

\newcommand{\exerciseheading}[2]{%
  \par\noindent
  {\large\bfseries #1\par}
  \vspace{0.1em}
  \noindent{\large #2\par}
  \vspace{1.5em}
}

\usepackage{amsmath}
\usepackage{amssymb}
\usepackage{amsthm}
\usepackage{mathtools}
\usepackage{aliascnt}

\usepackage[hidelinks,pdfusetitle]{hyperref}

\numberwithin{equation}{section}
\setlist{itemsep=0.25em,topsep=0.5em}

\theoremstyle{plain}
\newtheorem{theorem}{\notestheoremname}[section]
\newaliascnt{lemma}{theorem}
\newtheorem{lemma}[lemma]{\noteslemmaname}
\aliascntresetthe{lemma}
\newaliascnt{proposition}{theorem}
\newtheorem{proposition}[proposition]{\notespropositionname}
\aliascntresetthe{proposition}
\newaliascnt{corollary}{theorem}
\newtheorem{corollary}[corollary]{\notescorollaryname}
\aliascntresetthe{corollary}

\theoremstyle{definition}
\newaliascnt{definition}{theorem}
\newtheorem{definition}[definition]{\notesdefinitionname}
\aliascntresetthe{definition}
\newaliascnt{example}{theorem}
\newtheorem{example}[example]{\notesexamplename}
\aliascntresetthe{example}
\newaliascnt{problem}{theorem}
\newtheorem{problem}[problem]{\notesproblemname}
\aliascntresetthe{problem}
\newtheorem{exercise}{\notesexercisename}

\theoremstyle{remark}
\newaliascnt{remark}{theorem}
\newtheorem{remark}[remark]{\notesremarkname}
\aliascntresetthe{remark}

\newenvironment{solution}
{
\begin{proof}[\notessolutionname]}
  {
\end{proof}}

\usepackage[nameinlink,noabbrev]{cleveref}

\crefname{theorem}{theorem}{theorems}
\crefname{lemma}{lemma}{lemmas}
\crefname{proposition}{proposition}{propositions}
\crefname{corollary}{corollary}{corollaries}
\crefname{definition}{definition}{definitions}
\crefname{example}{example}{examples}
\crefname{problem}{problem}{problems}
\crefname{exercise}{exercise}{exercises}
\crefname{remark}{remark}{remarks}

\ifx\notelanguage\notesfinnishlanguage
\crefname{theorem}{lause}{lauseet}
\crefname{lemma}{lemma}{lemmat}
\crefname{proposition}{propositio}{propositiot}
\crefname{corollary}{seurauslause}{seurauslauseet}
\crefname{definition}{määritelmä}{määritelmät}
\crefname{example}{esimerkki}{esimerkit}
\crefname{problem}{tehtävä}{tehtävät}
\crefname{exercise}{tehtävä}{tehtävät}
\crefname{remark}{huomautus}{huomautukset}
\fi

\newcommand{\NN}{\mathbb{N}}
\newcommand{\NNone}{\mathbb{N}_1}
\newcommand{\ZZ}{\mathbb{Z}}
\newcommand{\QQ}{\mathbb{Q}}
\newcommand{\RR}{\mathbb{R}}
\newcommand{\CC}{\mathbb{C}}
\newcommand{\PP}{\mathbb{P}}

\DeclareMathOperator{\im}{im}
\DeclareMathOperator{\rank}{rank}
\DeclarePairedDelimiter{\abs}{\lvert}{\rvert}
\DeclarePairedDelimiter{\norm}{\lVert}{\rVert}
\newcommand{\set}[1]{\left\lbrace#1\right\rbrace}
\newcommand{\maxset}[1]{\max\set{#1}}
\newcommand{\minset}[1]{\min\set{#1}}
\newcommand{\powerset}[1]{\mathcal{P}\!\left(#1\right)}

\usepackage{forest}

\title{Johdatus logiikkaan I\\Harjoitus 6}
\author{}
\date{}

\begin{document}

\exerciseheading{Johdatus logiikkaan I}{Harjoituksen 6 ratkaisut}

\begin{exercise}
  Osoita \(\{\neg A\} \vdash A \to B\).
\end{exercise}

\begin{solution}
  \[
    \dfrac{
      \dfrac{
        \dfrac{[\neg B]^2\qquad
        \dfrac{[A]^1\qquad\neg A}{A\land\neg A}\;\land T}
        {\neg\neg B}\;\neg T,2
      }{B}\;\neg E
    }{A\to B}\;\to T,1
  \]
  \(\neg T\) hylkää käyttämättömän oletuksen \(\neg B\).
  Avoimeksi jää \(\neg A\).
\end{solution}

\begin{exercise}
  Osoita
  \[
    \{(A \land B) \to C\} \vdash (A \to C) \lor (B \to C).
  \]
\end{exercise}

\begin{solution}
  Merkitään \(D=(A\to C)\lor(B\to C)\) ja \(E=A\lor\neg A\).
  Johdetaan ensin \(E\) materiaalin esimerkin 7.6 tapaan:
  \[
    \mathcal L:\quad
    \dfrac{
      \dfrac{\dfrac{[A]^1}{E}\;\lor T\qquad[\neg E]^2}
      {E\land\neg E}\;\land T
    }{\neg A}\;\neg T,1
  \]
  \[
    \dfrac{
      \dfrac{
        \dfrac{
          \dfrac{
            \begin{array}{c}\mathcal L\\\neg A
          \end{array}}{E}\;\lor T
          \qquad[\neg E]^2
        }{E\land\neg E}\;\land T
      }{\neg\neg E}\;\neg T,2
    }{E}\;\neg E
  \]
  Tapauksessa \(A\):
  \[
    \mathcal P:\quad
    \dfrac{
      \dfrac{
        \dfrac{(A\land B)\to C\qquad
        \dfrac{[A]^3\qquad[B]^4}{A\land B}\;\land T}
        {C}\;\to E
      }{B\to C}\;\to T,4
    }{D}\;\lor T.
  \]
  Tapauksessa \(\neg A\):
  \[
    \mathcal Q:\quad
    \dfrac{
      \dfrac{
        \dfrac{
          \dfrac{[\neg C]^7\qquad
            \dfrac{[A]^6\qquad[\neg A]^5}
          {A\land\neg A}\;\land T}
          {\neg\neg C}\;\neg T,7
        }{C}\;\neg E
      }{A\to C}\;\to T,6
    }{D}\;\lor T.
  \]
  \(\neg T\) hylkää käyttämättömän oletuksen \(\neg C\).
  Yhdistetään tapaukset:
  \[
    \dfrac{
      E\qquad
      \begin{array}{c}\mathcal P\\D
      \end{array}\qquad
      \begin{array}{c}\mathcal Q\\D
      \end{array}
    }{D}\;\lor E,3,5.
  \]
  Avoimeksi jää \((A\land B)\to C\).
\end{solution}

\begin{exercise}
  Todista eheyslauseen kohdat 6 ja 7, eli osoita, että implikaation tuonti- ja
  eliminointisäännöt säilyttävät eheyden.
\end{exercise}

\begin{solution}
  \begin{enumerate}[label=\arabic*.,start=6]
    \item \emph{Implikaation tuonti.} Oletetaan induktio-oletuksena,
      että päättely \(P\), jonka johtopäätös on \(B\), on ehyt.
      Olkoon \(Q\) päättely
      \[
        Q:\quad
        \dfrac{
          \begin{array}{c}[A]\\P\\B
        \end{array}}{A\to B}\;\to T.
      \]
      Olkoon \(v\) totuusjakauma, joka toteuttaa \(Q\):n oletukset.
      Jos \(v(A)=0\), niin \(v(A\to B)=1\). Jos \(v(A)=1\),
      \(v\) toteuttaa myös päättelyssä \(Q\) hylätyt \(A\):n
      esiintymät ja siten kaikki \(P\):n oletukset. Induktio-oletuksen
      nojalla \(v(B)=1\), joten jälleen \(v(A\to B)=1\). Siis
      \(Q\) on ehyt.

    \item \emph{Implikaation eliminointi.} Oletetaan, että päättelyt
      \(P\), jonka johtopäätös on \(A\to B\), ja \(Q\), jonka
      johtopäätös on \(A\), ovat ehyitä. Olkoon \(R\) päättely
      \[
        R:\quad
        \dfrac{
          \begin{array}{c}P\\A\to B
          \end{array}
          \qquad
          \begin{array}{c}Q\\A
        \end{array}}{B}\;\to E.
      \]
      Jos \(v\) toteuttaa \(R\):n oletukset, se toteuttaa myös
      \(P\):n ja \(Q\):n oletukset. Induktio-oletuksen nojalla
      \(v(A\to B)=v(A)=1\), joten \(v(B)=1\). Siis \(R\) on ehyt.
  \end{enumerate}
\end{solution}

\begin{exercise}
  Osoita, ettei luonnollisella päättelyllä voi päätellä lausetta
  \((p_0 \lor p_1) \to p_2\) lauseesta \((p_0 \to p_2) \lor (p_1 \to p_2)\).
\end{exercise}

\begin{solution}
  Vastaesimerkki:
  \[
    v(p_0)=0,\qquad v(p_1)=1,\qquad v(p_2)=0.
  \]
  \[
    v\bigl((p_0\to p_2)\lor(p_1\to p_2)\bigr)=1,
    \qquad v\bigl((p_0\lor p_1)\to p_2\bigr)=0.
  \]
  Eheyden nojalla päättelyä ei ole.
\end{solution}

\begin{exercise}
  Voidaanko luonnollisella päättelyllä päätellä lause
  \((p_0 \to p_2) \land (p_1 \to p_2)\) lauseesta \((p_0 \land p_1) \to p_2\)?
\end{exercise}

\begin{solution}
  Ei voida. Vastaesimerkki:
  \[
    v(p_0)=0,\qquad v(p_1)=1,\qquad v(p_2)=0.
  \]
  \[
    v\bigl((p_0\land p_1)\to p_2\bigr)=1,
    \qquad v\bigl((p_0\to p_2)\land(p_1\to p_2)\bigr)=0.
  \]
  Eheyden nojalla päättelyä ei ole.
\end{solution}

\begin{exercise}
  Voidaanko luonnollisella päättelyllä päätellä lause
  \((p_0 \land p_1) \to p_2\) lauseesta
  \((p_0 \to p_2) \land (p_1 \to p_2)\)?
\end{exercise}

\begin{solution}
  Voidaan:
  \[
    \dfrac{
      \dfrac{
        \dfrac{(p_0\to p_2)\land(p_1\to p_2)}{p_0\to p_2}\;\land E
        \qquad
        \dfrac{[p_0\land p_1]^1}{p_0}\;\land E
      }{p_2}\;\to E
    }{(p_0\land p_1)\to p_2}\;\to T,1
  \]
  Avoimeksi jää \((p_0\to p_2)\land(p_1\to p_2)\).
\end{solution}

\begin{exercise}
  Etsi semanttisen puun avulla totuusjakauma, joka toteuttaa propositiolauseen
  \[
    (\neg(p_0 \land p_1) \to p_2) \land ((\neg p_0 \land \neg p_1)
    \to \neg p_2).
  \]
\end{exercise}

\begin{solution}
  Riittää jatkaa yhtä oksaa literaaleihin asti:
  \begin{center}
    \begin{forest}
      for tree={grow'=south, l sep=5mm, s sep=8mm,
      edge={-}, align=center}
      [$(\neg(p_0\land p_1)\to p_2)\land
        ((\neg p_0\land\neg p_1)\to\neg p_2)$
        [$\neg(p_0\land p_1)\to p_2$
          [$(\neg p_0\land\neg p_1)\to\neg p_2$
            [$\neg(\neg p_0\land\neg p_1)$
              [ei jatketa]
            ]
            [$\neg p_2$
              [$\neg\neg(p_0\land p_1)$
                [$p_0\land p_1$
                  [$p_0$
                    [$p_1$
                      [avoin]
                    ]
                  ]
                ]
              ]
              [$p_2$
                [$\times$]
              ]
            ]
          ]
        ]
      ]
    \end{forest}
  \end{center}
  Avoin oksa antaa toteuttavan jakauman
  \[
    v(p_0)=v(p_1)=1,\qquad v(p_2)=0.
  \]
\end{solution}

\begin{exercise}
  Osoita semanttisen puun avulla, että
  \[
    ((A \to C) \land (B \to C)) \to ((A \lor B) \to C)
  \]
  on tautologia.
\end{exercise}

\begin{solution}
  Lauseen negaation semanttinen puu:
  \begin{center}
    \begin{forest}
      for tree={grow'=south, l sep=4mm, s sep=9mm,
      edge={-}, align=center}
      [$\neg\bigl(((A\to C)\land(B\to C))\to((A\lor B)\to C)\bigr)$
        [$(A\to C)\land(B\to C)$
          [$\neg((A\lor B)\to C)$
            [$A\to C$
              [$B\to C$
                [$A\lor B$
                  [$\neg C$
                    [$A$
                      [$\neg A$
                        [$\times$]
                      ]
                      [$C$
                        [$\times$]
                      ]
                    ]
                    [$B$
                      [$\neg B$
                        [$\times$]
                      ]
                      [$C$
                        [$\times$]
                      ]
                    ]
                  ]
                ]
              ]
            ]
          ]
        ]
      ]
    \end{forest}
  \end{center}
  Kaikki oksat sulkeutuvat, joten lause on tautologia.
\end{solution}

\begin{exercise}
  Anna semanttinen todistus lauseelle \((A \to \neg B) \to \neg(A \land B)\).
\end{exercise}

\begin{solution}
  Lauseen negaation semanttinen puu:
  \begin{center}
    \begin{forest}
      for tree={grow'=south, l sep=4mm, s sep=12mm,
      edge={-}, align=center}
      [$\neg\bigl((A\to\neg B)\to\neg(A\land B)\bigr)$
        [$A\to\neg B$
          [$\neg\neg(A\land B)$
            [$A\land B$
              [$A$
                [$B$
                  [$\neg A$
                    [$\times$]
                  ]
                  [$\neg B$
                    [$\times$]
                  ]
                ]
              ]
            ]
          ]
        ]
      ]
    \end{forest}
  \end{center}
  Molemmat oksat sulkeutuvat, joten lause on tautologia.
\end{solution}

\begin{exercise}
  Anna semanttinen todistus lauseelle \(\neg(A \lor B)
  \leftrightarrow (\neg A \land \neg B)\).
\end{exercise}

\begin{solution}
  Lauseen negaation semanttinen puu:
  \begin{center}
    \begin{forest}
      for tree={grow'=south, l sep=4mm, s sep=8mm,
      edge={-}, align=center}
      [$\neg\bigl(\neg(A\lor B)\leftrightarrow
        (\neg A\land\neg B)\bigr)$
        [$\neg(A\lor B)$
          [$\neg(\neg A\land\neg B)$
            [$\neg A$
              [$\neg B$
                [$\neg\neg A$
                  [$A$
                    [$\times$]
                  ]
                ]
                [$\neg\neg B$
                  [$B$
                    [$\times$]
                  ]
                ]
              ]
            ]
          ]
        ]
        [$\neg\neg(A\lor B)$
          [$\neg A\land\neg B$
            [$A\lor B$
              [$\neg A$
                [$\neg B$
                  [$A$
                    [$\times$]
                  ]
                  [$B$
                    [$\times$]
                  ]
                ]
              ]
            ]
          ]
        ]
      ]
    \end{forest}
  \end{center}
  Kaikki oksat sulkeutuvat, joten ekvivalenssi on tautologia.
\end{solution}

\bigskip
\noindent\emph{Seuraavasta haastetehtävästä ei saa
  harjoituspisteitä. Se on tarkoitettu lisäpohdintatehtäväksi
aiheesta kiinnostuneille.}

\medskip
\noindent\textbf{Haastetehtävä.} Shefferin viiva on kaksipaikkainen
konnektiivi \(|\), jolle \(v(A|B) = 0\) jos ja
vain jos \(v(A) = v(B) = 1\). Anna \(A|B\):lle ja \(\neg(A|B)\):lle
säännöt semanttista puuta varten.
Riittävätkö säännöt ehyeeseen ja täydelliseen semanttiseen
todistukseen lauseille, joissa
esiintyy konnektiivina pelkkää Shefferin viivaa?

\end{document}
